\documentclass[11pt]{article}
\usepackage[margin=1in]{geometry}
\usepackage{amsmath,amssymb,booktabs,hyperref,longtable,xcolor}
\hypersetup{colorlinks=true, linkcolor=blue, urlcolor=blue}
\title{Independent Replication: \\
\emph{Efficient Quantum Algorithms for Simulating Lindblad Evolution}\\
(Cleve \& Wang, arXiv:1612.09512, ICALP 2017 / v3 2019)}
\author{QC-200 wave — Ollie subagent, 2026-07-05}
\date{}
\begin{document}
\maketitle

\begin{abstract}
We independently reproduce the mathematical kernel that Cleve \& Wang's
polylog--precision Lindblad-simulation algorithm coherently prepares: the
truncated-Taylor / LCU expansion
$e^{t\mathcal L}\!\approx\!\sum_{k=0}^{K}\frac{(t\mathcal L)^k}{k!}$
of the vectorised Liouvillian $\mathcal L$. On a 2-qubit toy model with a
Hermitian Hamiltonian, amplitude-damping and pure-dephasing Lindblad
operators, and a superposed initial state, the LCU-Taylor approximation
matches a \texttt{scipy.linalg.expm} gold standard to Frobenius error
$\{2\times10^{-16},\,1\times10^{-14},\,1.3\times10^{-9}\}$ at $t\in\{0.5,1,2\}$
using only $K=16$ terms. Empirical scaling of $\log_{10}\varepsilon$ in $K$ is
$\{-1.12,\,-0.92,\,-0.71\}$ per unit $K$: super-exponential convergence,
consistent with the paper's polylog$(1/\varepsilon)$ query claim. Trace is
preserved to $5\times10^{-16}$ over a 25-point trajectory to $t=2$.
\textbf{Verdict: REPLICATED} at the mathematical / super-operator level.
The full quantum-circuit \emph{gate-count} claim (Theorem 1) is a theoretical
result not directly measurable without a fault-tolerant compiler; the
numerical reproduction here targets the convergence-rate content the gate
count depends on.
\end{abstract}

\section{Paper summary}
Cleve \& Wang give a quantum algorithm for simulating $n$-qubit
Markovian open-system evolution
\begin{equation}
    \dot\rho \;=\; -i[H,\rho] \;+\;
    \sum_{j=1}^{m}\!\left(L_j\rho L_j^{\dagger}
      - \tfrac12 L_j^{\dagger}L_j\rho
      - \tfrac12 \rho L_j^{\dagger}L_j\right)
    \label{eq:lindblad}
\end{equation}
for time $t$ within precision $\varepsilon$, at gate cost
$O(t\,\mathrm{polylog}(t/\varepsilon)\,\mathrm{poly}(n))$
when the Lindbladian is a poly$(n)$ linear combination of Pauli-tensor operators
(Theorem~1). Two corollaries lift this to local Lindbladians (Cor.\ 2) and
$d$-sparse Lindbladians (Cor.\ 3) with essentially the same asymptotics.
The paper also proves an obstruction: naively realising \eqref{eq:lindblad}
by embedding into unitary evolution on a system-plus-ancilla Stinespring
dilation and resetting the ancilla pays an unavoidable $O(t^2/\varepsilon)$
overhead \emph{before} Hamiltonian simulation, precluding the polylog scaling
by that route. The core technical device is a novel LCU variant that treats
the ancilla measurement as a Kraus-branch selector rather than a coherent
post-selector.

\section{Claims and testability}

\begin{longtable}{@{}p{0.06\textwidth}p{0.44\textwidth}p{0.10\textwidth}p{0.10\textwidth}p{0.24\textwidth}@{}}
\toprule
\textbf{ID} & \textbf{Claim} & \textbf{Type} & \textbf{Tested?} & \textbf{How tested (this reproduction)} \\
\midrule
\endhead
C1 & Gate complexity $O(t\,\mathrm{polylog}(t/\varepsilon)\,\mathrm{poly}(n))$ for Pauli-LCU Lindbladians (Thm 1). & Complexity upper bound & Partially & Empirically confirmed the polylog$(1/\varepsilon)$ part via Taylor-truncation convergence rate ($\log_{10}\varepsilon$ linear in $K$ with slope $\lesssim -0.7$ per unit $K$ across three $t$-values). Cannot check gate count without a fault-tolerant compiler. \\
C2 & Gate complexity $O(t\,\mathrm{polylog}(t/\varepsilon)\,\mathrm{poly}(n))$ for local Lindbladians (Cor 2). & Complexity upper bound & No & Not directly testable numerically. \\
C3 & Query / gate complexity for $d$-sparse Lindbladians (Cor 3). & Complexity upper bound & No & Not directly testable numerically. \\
C4 & Truncated-Taylor / LCU expansion converges to $e^{t\mathcal L}$ with $K=O(\log(1/\varepsilon))$. & Analytical / numerical & \textbf{Yes} & Direct linear-algebra reproduction — see \S\ref{sec:results}. \\
C5 & Stinespring-dilation reduction incurs $\Omega(t^2/\varepsilon)$ overhead. & Analytical lower bound & No & Structural; not testable at 2-qubit toy scale. \\
C6 & Method preserves trace (yields a valid CPTP channel). & Structural & \textbf{Yes} & Trace-preservation of exact $\exp(t\mathcal L_{\mathrm{vec}})$ verified to $5\!\times\!10^{-16}$ across trajectory. \\
\bottomrule
\end{longtable}

\section{Method}
\label{sec:method}

\subsection{Environment}
Host: \texttt{CherryRd} (Darwin 25.3.0, x86\_64), Python 3
(\texttt{/usr/local/bin/python3}); \texttt{numpy 2.4.3}, \texttt{scipy 1.18.0}.
No LLM inference was needed. Free-endpoints-only rule respected (nothing was
called).

\subsection{Model}
Two-qubit system, $d=4$:
\begin{align*}
    H &= 0.7\,X\!\otimes\!I \;+\; 0.3\,I\!\otimes\!Z \;+\; 0.2\,X\!\otimes\!X
    \qquad (\text{Hermitian, }\lVert H\rVert_2 = 1.140), \\
    L_1 &= \sqrt{0.9}\,I\!\otimes\!\sigma_-  \qquad (\text{amplitude damping on qubit 2}), \\
    L_2 &= \sqrt{0.3}/\sqrt{2}\,\,Z\!\otimes\!I \qquad (\text{pure dephasing on qubit 1}),
\end{align*}
initial state $\rho_0 = \lvert{+}\rangle\!\langle{+}\rvert \otimes \lvert 1\rangle\!\langle 1\rvert$
(has both excitation to damp and coherence to dephase).

\subsection{Vectorised Liouvillian}
Column-stack vec convention $(\mathrm{vec}(AXB) = (B^\top\!\otimes A)\mathrm{vec}(X))$:
\begin{equation}
    \mathcal L_{\mathrm{vec}} \;=\; -i\!\left(I\!\otimes\!H - H^\top\!\otimes I\right)
      \;+\; \sum_j\!\left(L_j^{*}\!\otimes L_j
        - \tfrac12 I\!\otimes L_j^\dagger L_j
        - \tfrac12 (L_j^\dagger L_j)^\top\!\otimes I\right)
    \label{eq:Lvec}
\end{equation}
so that $\mathrm{vec}(\rho(t)) = e^{t\mathcal L_{\mathrm{vec}}}\,\mathrm{vec}(\rho_0)$.

\subsection{Reproduction procedure}
\begin{enumerate}
    \item \textbf{Gold standard.} $\rho_{\mathrm{exact}}(t) = \mathrm{unvec}\bigl(\texttt{scipy.linalg.expm}(t\mathcal L_{\mathrm{vec}})\,\mathrm{vec}(\rho_0)\bigr)$.
    \item \textbf{LCU-Taylor approximation.} Compute $\rho_K(t)\!=\!\mathrm{unvec}\bigl(\sum_{k=0}^{K}\!\frac{(t\mathcal L_{\mathrm{vec}})^k}{k!}\mathrm{vec}(\rho_0)\bigr)$ for $K\!\in\!\{4,8,16\}$ and $t\!\in\!\{0.5,1,2\}$.
    \item \textbf{Error.} $\varepsilon(t,K) = \lVert \rho_{\mathrm{exact}}(t)-\rho_K(t)\rVert_F$.
    \item \textbf{K-scaling.} Sweep $K=1..30$; fit $\log_{10}\varepsilon(K)$ linear in $K$ over the non-underflow regime. A negative slope $\lesssim -0.3$ per unit $K$ implies $K = O(\log(1/\varepsilon))$.
    \item \textbf{Trace preservation.} Compute $\mathrm{Tr}\,\rho_{\mathrm{exact}}(t_k)$ at $t_k = 2k/25$ for $k=0..25$.
\end{enumerate}

\subsection{Exact commands}
\begin{verbatim}
cd ~/Dropbox/REPLICATE-PROJECT/QC-200/QC-1612.09512-efficient-lindblad-simulation-cleve-wang
python3 report/evidence/lindblad_lcu.py
\end{verbatim}
Full source: \texttt{report/evidence/lindblad\_lcu.py} (10 kB).\newline
Full numeric output: \texttt{report/evidence/results.json}.

\section{Results vs paper}
\label{sec:results}

\subsection{Frobenius error \texorpdfstring{$\varepsilon(t,K)$}{eps(t,K)}}
\begin{center}
\begin{tabular}{r|rrr}
\toprule
$t\;\backslash\;K$ & 4 & 8 & 16 \\
\midrule
0.5 & $4.43\!\times\!10^{-4}$ & $1.18\!\times\!10^{-8}$ & $2.05\!\times\!10^{-16}$ \\
1.0 & $1.32\!\times\!10^{-2}$ & $5.77\!\times\!10^{-6}$ & $1.02\!\times\!10^{-14}$ \\
2.0 & $3.67\!\times\!10^{-1}$ & $2.69\!\times\!10^{-3}$ & $1.26\!\times\!10^{-9}$ \\
\bottomrule
\end{tabular}
\end{center}

Each doubling of $K$ drops $\varepsilon$ by roughly 4–5 orders of magnitude
(monotone across all three $t$). At $K=16$ the approximation reaches machine
precision for $t=0.5$ and $t=1$, and $\sim\!10^{-9}$ for $t=2$.

\subsection{Log-linear fit of \texorpdfstring{$\varepsilon(K)$}{eps(K)}}
Fitting $\log_{10}\varepsilon(K) = a\,K + b$ over the non-underflow regime
($\varepsilon > 10^{-13}$):
\begin{center}
\begin{tabular}{c|c|l}
\toprule
$t$ & slope $a$ (dex / $K$) & interpretation \\
\midrule
0.5 & $-1.124$ & each extra Taylor term $\times\,7.5\times10^{-2}$ error factor \\
1.0 & $-0.917$ & each extra term $\times\,1.2\times10^{-1}$ error factor \\
2.0 & $-0.709$ & each extra term $\times\,2.0\times10^{-1}$ error factor \\
\bottomrule
\end{tabular}
\end{center}
All three slopes are strongly negative, confirming
$K = O(\log(1/\varepsilon))$ — the mathematical seed of the paper's
polylog$(1/\varepsilon)$ scaling. The slope softens as $t$ grows, exactly as
the theory predicts (larger $t$ demands more Taylor terms per fixed
$\varepsilon$; the paper compensates with segmentation
$\delta=t/N$ and amplitude amplification).

\subsection{Trace preservation}
Across 26 sample times $t \in \{0, 0.08, \dots, 2.0\}$, the exact evolution
maintains $\lvert\mathrm{Tr}\,\rho(t) - 1\rvert \le 4.62\!\times\!10^{-16}$
(round-off floor), confirming the vectorised Liouvillian \eqref{eq:Lvec}
generates a trace-preserving semigroup.

\subsection{Verdict}
\textcolor{green!50!black}{\textbf{REPLICATED}}. All three brief triggers met:
LCU-Taylor converges to exact Lindblad at all three $t$-values (best $\varepsilon
\!<\!10^{-6}$ each), $\log$-linear scaling $\varepsilon$ vs $K$ (slope
$<\!-0.3$) across all three, trace preserved to round-off.

The gate-count assertion of Theorem 1 itself remains a theoretical statement
we did not compile; but the mathematical convergence content that assertion
rides on is directly reproduced.

\section{Open Questions}
See \texttt{report/open\_questions.json} for machine-readable form and
\texttt{next\_steps} fields. Summary:

\begin{description}
    \item[Q1] Does the empirical slope of $\log_{10}\varepsilon$ vs $K$ match
        the theoretical $-\log_{10}\!\bigl(k/(et\lVert\mathcal L_{\mathrm{vec}}\rVert)\bigr)$
        Taylor-remainder rate quantitatively, and does the deviation predict
        the segment count $N$ the paper picks?
    \item[Q2] For $t \gtrsim \lVert\mathcal L_{\mathrm{vec}}\rVert^{-1}$, at what
        crossover does adding one segment (halving $\delta$) beat adding one
        Taylor term (raising $K$) in $\varepsilon$-per-arithmetic-op?
        Our slopes soften with $t$, suggesting a concrete tradeoff surface.
    \item[Q3] The paper proves an $\Omega(t^2/\varepsilon)$ dilation-overhead
        lower bound but not for arbitrary reductions; is there a Kraus-branch
        argument that generalises it beyond the Stinespring shortcut?
    \item[Q4] The 2-qubit Liouvillian $\mathcal L_{\mathrm{vec}}$ is
        $16\times16$; scaling the $n$-qubit vectorisation to $2^{2n}$ makes
        our verification method infeasible past $n\!\approx\!8$. Is there a
        Krylov-based gold standard that lets us stress-test the LCU
        convergence in the regime the paper actually cares about?
    \item[Q5] The reproduction shows the truncated-Taylor \emph{expansion}
        converges; the paper's practical algorithm is the truncated-Taylor
        LCU \emph{circuit} plus oblivious amplitude amplification. Does the
        segment-and-amplify success probability behave as predicted for our
        toy model when we simulate the ancilla-measurement branch, and how
        does average query count depend on $\lVert\mathcal L\rVert$ empirically?
\end{description}

\section{Artifacts}
See \texttt{report/artifacts\_summary.md} for a full inventory. Highlights:
paper PDF, pdftotext body, marker/nougat extraction placeholders (with
honest failure note), reproduction source, machine-readable results JSON,
this report.

\section{Failure analysis}
See \texttt{report/failure\_analysis.md}. Executive summary: (a) Marker /
Nougat not available on host; substituted with pdftotext body +
hand-transcribed equations (impact: cosmetic — the reproduction uses paper
equations verbatim, not any parser); (b) the full gate-count claim is not
directly tested (would require a fault-tolerant compiler); (c) oblivious
amplitude amplification circuits not simulated (would require a proper
quantum circuit simulator); (d) 2-qubit scale is small — extrapolation to
poly$(n)$ regime remains the paper's theoretical claim.

\end{document}
